% Section I. Target: ~1 page (6-page budget, A3 2026-11-18).
% Evidence: research/comparison.csv, research/limitations-evidence.md,
% research/peru-energy-stats.md (abstract-level; verified 2026-10-28).
\section{Introducción}
\label{sec:introduction}

% A. Problem context.
En el Perú, la producción de electricidad alcanzó \SI{63565}{GWh} en 2024, \SI{2,4}{\percent} más que el año anterior~\cite{minem2025produccion}, y en el primer trimestre de ese año el \SI{96,0}{\percent} de la población accedía a la red pública~\cite{inei2024acceso}.
Con esa cobertura, el reto se desplaza del acceso hacia el uso informado y eficiente de la energía.
Sin embargo, la supervisión del consumo suele basarse en la facturación mensual, que no revela variaciones de demanda, picos ni patrones horarios.
Los sistemas de Internet de las Cosas (IoT) miden variables eléctricas con mayor frecuencia y las analizan de forma continua~\cite{rind2023perspective}, pero su utilidad depende de la exactitud de las mediciones, la calidad de los datos y la latencia de transmisión.

% B. Related work.
La literatura reciente ofrece numerosos medidores IoT de bajo costo.
Sousa \textit{et~al.}~\cite{sousa2023smartmeter} desarrollaron un medidor inalámbrico con ESP32 y medición de calidad de energía, y Garcés \textit{et~al.}~\cite{garces2025iotmeter} uno con sensores no invasivos y calibración manual de señales.
Munoz \textit{et~al.}~\cite{munoz2024smartmeter} construyeron un medidor trifásico de menos de \SI{80}{USD} con error absoluto menor al \SI{1}{\percent}, mientras que Rohman \textit{et~al.}~\cite{rohman2025esp32} emplearon un módulo integrado PZEM-016 con un servidor web local.
En cuanto a la validación, Kizonde \textit{et~al.}~\cite{kizonde2023monitoringnode} reportaron baja linealidad en cargas pequeñas y sesgo en cargas grandes frente a una unidad de medición de potencia.
Cárdenas Magaña \textit{et~al.}~\cite{cardenas2026residential} validaron un sistema ESP32 con PZEM-004T frente a un instrumento calibrado con cargas resistivas, inductivas y electrónicas, con errores absolutos medios de \SI{0,31}{V}, \SI{0,008}{A} y \SI{1,16}{W}.
Marin \textit{et~al.}~\cite{marin2026rmsbenchmark} contrastaron mediciones RMS de nodos con sensores de efecto Hall frente a un analizador de calidad de energía.
Respecto a la comunicación, Khan y Pirak~\cite{khan2021mqtt} compararon MQTT y CoAP en un medidor NB-IoT, Nizam \textit{et~al.}~\cite{nizam2024realtime} integraron Modbus, ESP-NOW y MQTT con actualizaciones cada \SI{5}{s}, y De Sousa \textit{et~al.}~\cite{desousa2026smartmeter} usaron ESP32-S3 con MQTT, InfluxDB y Grafana, con un error medio de \SI{4,56}{\percent} en corriente y latencias de \SIrange{203,63}{290,86}{ms} en el servidor.
En el análisis de datos, Oprea \textit{et~al.}~\cite{oprea2021anomaly} detectaron anomalías con una exactitud del \SI{90}{\percent} y Zheng \textit{et~al.}~\cite{zheng2023thresholds} propusieron umbrales dinámicos para picos anómalos.

% C. Limitations and gap (narrowed to what the evidence supports, AD-8).
Estos trabajos presentan limitaciones recurrentes.
Primero, la validación metrológica es heterogénea: algunos no identifican el instrumento de referencia o evalúan pocas cargas~\cite{rohman2025esp32}, y otros reportan el error de una sola magnitud~\cite{desousa2026smartmeter}.
Segundo, la latencia extremo a extremo rara vez se mide; entre los trabajos revisados, solo~\cite{desousa2026smartmeter} la reporta.
Tercero, aunque hay sistemas que almacenan datos sin conexión~\cite{kizonde2023monitoringnode}, no se reporta una evaluación controlada de la recuperación tras una interrupción de la comunicación.
Cuarto, los métodos de análisis se evalúan sobre conjuntos de datos ajenos al hardware de medición~\cite{oprea2021anomaly,zheng2023thresholds}.
No se identificó un prototipo de bajo costo que evalúe conjuntamente la exactitud frente a una referencia, la latencia, la recuperación ante interrupciones y los patrones de consumo.

% D. Objective and contributions (aligned with spec 001; no results stated).
El objetivo de este trabajo es diseñar, implementar y evaluar experimentalmente un sistema IoT para el monitoreo y análisis en tiempo real flexible del consumo eléctrico de un circuito monofásico controlado.
Las contribuciones son:
\begin{enumerate}
  \item un nodo de bajo costo, construido con componentes disponibles, que mide tensión y corriente con corrección de desfase entre canales y calcula localmente la potencia activa y la energía;
  \item una cadena de entrega MQTT con confirmación, cola local y almacenamiento idempotente en una base de series temporales;
  \item un análisis del consumo por hora, día y semana, con alertas por umbral que descartan los transitorios de arranque;
  \item una evaluación experimental de la exactitud frente a un instrumento de referencia, de la latencia y de la recuperación tras interrupciones controladas de la red.
\end{enumerate}
La Sección~\ref{sec:methodology} describe la metodología, la Sección~\ref{sec:results} los resultados y la Sección~\ref{sec:conclusions} las conclusiones.
